% docs/texte/10_engineering/version_control.tex -- chapter "Version Control" (Interferosim)
% Version: 0.1.3 (2026-11-03) -- repository public, MIT licence
%   0.1.2: workflow diagram in 16.3
%   0.1.1: complete command-line example in 16.3
% ================================================================

\index{version control}\index{git}\index{GitHub}
\section{Repository}

\begin{center}
\begin{tabular}{@{}ll@{}}
\toprule
Remote repository & \url{https://github.com/eschubi/interferosim} (public) \\
Default branch & \texttt{main} \\
Local working copy & project folder \texttt{\ldots\textbackslash sem11\textbackslash open-source\textbackslash interferosim} \\
Tools & GitHub Desktop or the \texttt{git} command line \\
\bottomrule
\end{tabular}
\end{center}

The repository is public; the code is released under the MIT License
(\texttt{LICENSE}).

\section{What is under version control}

The rules are in \texttt{.gitignore} and \texttt{.gitattributes} in the
project root.
\begin{itemize}
  \item \textbf{Included:} source code (\texttt{src/}), tests,
        notebooks, configuration files (\texttt{configs/}), the \LaTeX{}
        sources of the documentation including the local style copy
        (\texttt{docs/style/}), figures made in the project, README,
        CHANGELOG, LICENSE, CITATION.cff and the set-up scripts.
  \item \textbf{Excluded:} the virtual environment \texttt{.venv},
        Python caches, Jupyter checkpoints, \LaTeX{} auxiliary files,
        built PDFs (\texttt{docs/*.pdf}, they are generated from the
        sources), archives (\texttt{*.zip}, \texttt{*.tar.gz}) and all
        files in \texttt{docs/images/} except those made in the project.
        Third-party images stay local until their licences are clarified;
        a new project figure must be released explicitly in
        \texttt{.gitignore} (line \texttt{!docs/images/<name>}).
  \item \textbf{Line endings:} \texttt{.gitattributes} stores text files
        with LF; \texttt{.bat} and \texttt{.ps1} are checked out with CRLF,
        \texttt{.sh} with LF, PNG and PDF files are treated as binary.
\end{itemize}

\section{Daily work with GitHub Desktop}

Figure~\ref{fig:git-workflow} shows the course of a working session; the
steps are the same in GitHub Desktop and on the command line.

\begin{figure}[htbp]
  \centering
  % docs/texte/10_engineering/figures/git_workflow.tex -- workflow diagram of a working session (TikZ)
% Version: 0.1.0 (2026-10-12)
% Flow of chapter "Version Control", section "Daily work": update, work,
% test, review, commit, push; optional release with version, tag and
% push --tags. Left column: local working copy; right: GitHub.
\begin{tikzpicture}[
    node distance=5mm and 9mm,
    >=Stealth, font=\footnotesize\sffamily,
    wfstep/.style={draw, thick, rounded corners=2pt, fill=blue!6,
                 minimum width=46mm, minimum height=7mm, align=center},
    cmd/.style={font=\scriptsize\ttfamily, text=black!70},
    dec/.style={draw, thick, diamond, aspect=2.4, fill=orange!12,
                inner sep=1pt, align=center},
    term/.style={draw, thick, rounded corners=3.5mm, fill=gray!15,
                 minimum width=30mm, minimum height=7mm},
    remote/.style={draw, thick, fill=green!8, minimum width=30mm,
                   minimum height=16mm, align=center},
    flow/.style={->, thick},
  ]
  % ---- main flow (local) ---------------------------------------------
  \node[term] (start) {start of session};
  \node[wfstep, below=of start] (act)  {activate environment\\[-1pt]
        {\scriptsize\ttfamily .venv\textbackslash Scripts\textbackslash Activate.ps1}};
  \node[wfstep, below=of act]   (pull) {update working copy\\[-1pt]{\scriptsize\ttfamily git pull}};
  \node[wfstep, below=of pull]  (work) {work: edit code, notebooks, docs};
  \node[wfstep, below=of work]  (test) {run the tests\\[-1pt]{\scriptsize\ttfamily pytest -q}};
  \node[dec,  below=of test]  (ok)   {all tests\\pass?};
  \node[wfstep, below=of ok]    (rev)  {review changes\\[-1pt]{\scriptsize\ttfamily git status\quad git diff}};
  \node[wfstep, below=of rev]   (add)  {select files\\[-1pt]{\scriptsize\ttfamily git add <files>}};
  \node[wfstep, below=of add]   (com)  {commit\\[-1pt]{\scriptsize\ttfamily git commit -m "..."}};
  \node[dec,  below=of com]   (rel)  {release?};
  \node[wfstep, below=of rel]   (push) {upload\\[-1pt]{\scriptsize\ttfamily git push}};
  \node[term, below=22mm of push]  (end)  {end of session};

  % ---- remote: GitHub, right of "update working copy" --------------------
  \node[remote, right=32mm of pull] (gh)
        {\textbf{GitHub}\\[1pt]{\scriptsize\ttfamily eschubi/interferosim}\\[-1pt]
         {\scriptsize\ttfamily origin/main}};

  % ---- release branch, below GitHub in the right column -------------------
  \node[wfstep, fill=yellow!12, minimum width=44mm] (ver) at (rel -| gh)
        {raise version, CHANGELOG\\[-1pt]
         {\scriptsize\ttfamily pyproject.toml, \_\_init\_\_.py}};
  \node[wfstep, below=of ver, fill=yellow!12, minimum width=44mm] (rcom)
        {release commit and tag\\[-1pt]{\scriptsize\ttfamily git commit\quad git tag v0.3.0}};
  \node[wfstep, below=of rcom, fill=yellow!12, minimum width=44mm] (rpush)
        {upload with tags\\[-1pt]{\scriptsize\ttfamily git push\quad git push --tags}};

  % ---- arrows ----------------------------------------------------------
  \draw[flow] (start) -- (act);
  \draw[flow] (act)   -- (pull);
  \draw[flow] (pull)  -- (work);
  \draw[flow] (work)  -- (test);
  \draw[flow] (test)  -- (ok);
  \draw[flow] (ok)    -- node[right]{yes} (rev);
  \draw[flow] (ok.west) -- node[above]{no} ++(-12mm,0) |- node[pos=0.25,left]{fix} (work.west);
  \draw[flow] (rev)   -- (add);
  \draw[flow] (add)   -- (com);
  \draw[flow] (com)   -- (rel);
  \draw[flow] (rel)   -- node[right]{no} (push);
  \draw[flow] (rel)   -- node[above]{yes} (ver);
  \draw[flow] (ver)   -- (rcom);
  \draw[flow] (rcom)  -- (rpush);
  \draw[flow] (push)  -- (end);
  \draw[flow] (rpush.south) |- (end.east);

  % ---- data exchange with GitHub (dashed) ----------------------------------
  \coordinate (corr) at ($(pull.east)!0.45!(gh.west)$);   % corridor between the columns
  \draw[->, thick, dashed, green!45!black] ($(gh.west)+(0,3mm)$) -- node[above, font=\scriptsize]{pull}
        ($(pull.east)+(0,3mm)$);
  \draw[->, thick, dashed, green!45!black] (push.east) -- (push.east -| corr)
        node[pos=0.5, below, font=\scriptsize]{push}
        |- ($(gh.west)+(0,-3mm)$);
  \draw[->, thick, dashed, green!45!black] (rpush.east) -- ++(8mm,0)
        node[right, font=\scriptsize, align=left]{push\\tags} |- (gh.east);

  % columns
  \node[font=\scriptsize\itshape, text=black!60, above=1mm of start] {local working copy};
  \node[font=\scriptsize\itshape, text=black!60, above=1mm of gh] {remote};
\end{tikzpicture}

  \caption{Workflow of a working session: update, work, test, review,
           commit and upload; for a release additionally version,
           CHANGELOG and tag. Dashed: data exchange with GitHub.}
  \label{fig:git-workflow}
\end{figure}

\begin{enumerate}
  \item \textbf{Once:} \emph{File $\rightarrow$ Add local repository}
        and select the project folder. The repository is already
        connected to \texttt{origin} (GitHub).
  \item \textbf{Before starting:} \emph{Fetch origin} and, if offered,
        \emph{Pull origin}, so that the working copy is up to date.
  \item \textbf{Work}, then run the tests (chapter~\ref{ch:tests}).
  \item \textbf{Commit:} in the \emph{Changes} tab select only the files
        that belong to the change, write a summary and a description
        (see below), \emph{Commit to main}.
  \item \textbf{Push:} \emph{Push origin}.
\end{enumerate}

\paragraph*{Command-line equivalents} (in the project folder):
\begin{lstlisting}[basicstyle=\ttfamily\footnotesize, frame=single]
git status                    # what has changed
git pull                      # update the working copy
git add <files>               # select files for the commit
git commit -m "Summary"       # commit (opens an editor without -m)
git push                      # upload to GitHub
git log --oneline             # history
\end{lstlisting}

\paragraph*{Complete workflow on the command line.} A full working
session in PowerShell, from updating the working copy to the upload, here
for a change in the sky package:
\begin{lstlisting}[basicstyle=\ttfamily\footnotesize, frame=single]
# 1. go to the project and activate the project environment
cd C:\Users\erich\Dropbox\work-astro\Astro\Astronomie\sem11\open-source\interferosim
.venv\Scripts\Activate.ps1

# 2. get the current state from GitHub
git pull

# 3. work: edit files (editor, Jupyter, ...)

# 4. run the tests -- only continue if all pass
pytest -q

# 5. review the changes
git status                         # changed / new files
git diff                           # line-by-line changes

# 6. select the files that belong to this change
git add src/interferosim/sky/components.py tests/test_sky.py

# 7. commit: summary line, empty line, description
git commit -m "sky: add ring component" -m "RingSource with analytic visibility; tests added."

# 8. upload
git push

# 9. check
git log --oneline -3
git status                         # "Your branch is up to date with 'origin/main'"
\end{lstlisting}

\paragraph*{Release.} For a new version the same session ends with a
version change and a tag:
\begin{lstlisting}[basicstyle=\ttfamily\footnotesize, frame=single]
# version in pyproject.toml and src/interferosim/__init__.py raised to 0.3.0,
# entry in CHANGELOG.md written, tests passed
git add pyproject.toml src/interferosim/__init__.py CHANGELOG.md
git commit -m "Interferosim 0.3.0: <short summary>"
git tag v0.3.0
git push
git push --tags
\end{lstlisting}

\section{Commits and releases}

\begin{itemize}
  \item One commit per self-contained change (a module, a chapter, a
        fix). Files of unfinished work, test output and scratch notebooks
        (\texttt{Untitled.ipynb}) are not committed.
  \item Commit message: first line a short summary in English, e.\,g.\
        \texttt{sky: add DiscSource}; for a release the version number,
        e.\,g.\ \texttt{Interferosim 0.2.1: ...}. Further lines explain
        what and why.
  \item Every release increases the version in \texttt{pyproject.toml}
        and \texttt{src/interferosim/\_\_init\_\_.py} (chapter
        \enquote{Development Guidelines}) and gets an entry in
        \texttt{CHANGELOG.md}.
  \item Releases are tagged \texttt{v<major>.<minor>.<patch>}
        (GitHub Desktop: \emph{History}, right-click on the commit,
        \emph{Create Tag}; command line: \texttt{git tag v0.2.1} and
        \texttt{git push --tags}).
\end{itemize}

\section{Git and Dropbox}

The working copy lies in a Dropbox folder, so the \texttt{.git} folder is
synchronised as well. This is harmless as long as only one computer
works on the repository at a time. If several computers are used, each
should either wait until Dropbox has finished synchronising before
running git, or use its own clone outside Dropbox (\emph{File
$\rightarrow$ Clone repository} in GitHub Desktop) and exchange changes
only via GitHub. The virtual environment can be excluded from Dropbox
with \texttt{setup\_venv.bat -DropboxIgnore}.
